\chapter{Single-Species Transport}
\def\shorttitle{Single-Species Transport} %

\section{PHT3D Exercise , Part 1: Single-Species 2-D Transport Simulation}


\subsection{Groundwater Flow}


\begin{figure}[tb]
\centering
\includegraphics[height=0.9\textwidth, angle = 270]{mt3dms_ex1_3dview.eps}
\caption[Illustration of the conceptual model ]{Illustration of the conceptual model}
\label{fig:single_spec_fig1}
\end{figure}

\subsection{Solute Transport}

The 

\begin{table}[hb]
\caption{Flow and solute transport parameters used for PHT3D Exercise 7}
\begin{center}
\begin{tabular}{lc} 
\toprule
Parameter                                            &         Value \\
\midrule
Flow simulation type                                 &        steady state \\
Simulation time (days)                               &  21000           \\
Time step (days)                               &  60           \\
Model \  extent\ column\ direction\ $ (m)$           &  400            \\
Model \ top\ elevation\ $(m)$                        &   16            \\
Model \ bottom\ elevation\ $(m)$                        &   0            \\
Horizontal hydraulic conductivity\ $(m/day)$         &   110     \\ 
Porosity, $\theta$                                   &           0.39  \\
Piezometric head downstream boundary\ $(m) $         &  14            \\ 
Longitudinal dispersivity, $\alpha_L$   $(m)$             &           1.0 \\
Horizontal transverse dispersivity, $\alpha_{TH}$    &           0.1 \\
Vertical transverse dispersivity,   $\alpha_{TV}$    &           0.01 \\
Effective diffusion coefficient, D$^{\star}$          &             0  \\
\bottomrule

\end{tabular}
\label{single_spec_table_1}
\end{center}
\end{table}

\subsection{Setting up the flow model}

To set up the flow model proceed as follows:

\begin{itemize}

\item In iPHT3D, create a new model and save it in a new and separate directory (folder).

\item Select Parameters $\rightarrow$ 1.Model $\rightarrow$ Model to specify the model dimension ('Xsection')
and type ('free')
\item Select Parameters $\rightarrow$ 1.Model $\rightarrow$ Domain to specify
the domain as described in Table \ref{single_spec_table_1}. For the first simulation select a spacing of 20 m in x direction
and 1 m in y direction

\item Select Parameters $\rightarrow$ Time $\rightarrow$ and set the total Simulation Time to 21000 days. Then define the step size to 60 days

\item Set the boundary conditions for heads by selecting Spatial Attributes $\rightarrow$ Modflow $\rightarrow$ bas.3 boundary conditions corresponding to
IBOUND (Modflow). Use the zone tool to create a zone at the outflow from 0 to 13.9 m

\item to fix heads at the boundaries, use Spatial Attributes $\rightarrow$ Modflow $\rightarrow$ bas.5 Initial heads. provide a value of 16 in the dialog box.
this value will be used everywhere as a starting head and will be used as a BC at places where the BC where fixed (-1) previously. In order to 
set the outflow boundary, draw a line at the same place as the outflow BC and attribute a value of 14 m.

\item the inlet is done by a well between 0 and 15 m height with an total injection rate of 2.16 $m^3/d$

\item select Parameters $\rightarrow$ 2.Flow $\rightarrow$ Parameters and in this dialog set several parametters. 
In lpf.2 select convertible for the type of layer. Use lpf.8 to set the hydraulic conductivity.

\item to fix recharge, use Spatial Attributes $\rightarrow$ Modflow $\rightarrow$ rch.2 recharge. create a zone 
on the top cell with a value of 0.0014 in the dialog box.

\item for porosity value, Select Spatial Attributes $\rightarrow$ MT3DMS $\rightarrow$ btn.11 and set the value to 0.39. 

\end{itemize}
save the model changes.

\subsection{Running MODFLOW}

You have now completed to set up the flow model and you are ready to run MODFLOW.
This step will provide MT3DMS with the flow field that is used to compute 
advective-dispersive transport.

\begin{itemize}

\item Select Parameters $\rightarrow$ 2.Flow $\rightarrow$ Write

\item Select Parameters $\rightarrow$ 2.Flow $\rightarrow$ Run...

\item When the calculation has finished, a window appears, verify that 'NORMAL TERMINATION OF MODFLOW' is written.

\end{itemize}

\subsection{Visualising hydraulic heads}

\begin{itemize}

\item Visualise results by selecting Results $\rightarrow$ 2.Flow $\rightarrow$ Head. The time step can be changed
in Results $\rightarrow$ 1.Model $\rightarrow$ Time.

\end{itemize}

\subsection{Setting up the transport model}

to be copied from previous models

\begin{itemize}

\item Define the observation well Spatial Attributes $\rightarrow$ OBSERVATION $\rightarrow$ obs.1 observations
use the zone tool to draw a point at any location

\item visualize the head contours in Results $\rightarrow$ 3.Transport $\rightarrow$ Tracer. You can change the Time above 
to see another time. It is possible, after selecting one time to use the wheel mouse to stream along time.

\item Visualise the breakthrough curve by selecting Results $\rightarrow$ 5.Observation $\rightarrow$ Breakthrough, then below, 
select the name of the point you just create. A dialog appears that allows you to select value and a breakthrough curves appears.
The curve can be exported to be used in a spreadsheet.

\end{itemize}

\subsection{Chemistry}

in this section we will define the initial composition and solutions used as source with the following steps :
\item import the database
\item the phase in presence everywhere (background) is goethite at 0.1 mol/L. FeS(ppt) must also be ticked but with a null 
concentration (only precipitaiton allowed)
\item As previously the initial solution composition are entered in the Solution Tab, we need the background and two solutions.
the species needed to be activated are Cl,C(4), Fe(2) and Fe(3), Na, O(0), S(-2) and S(6), Orgc, Zn and like always pH and pe. 
Concentrations are given in table ??
\item the solutions are entered in the ph.4 source/sink line at inflow cells as zones as follows :
 from 14 to 16 m : solution 2, from 3 to 14 m : solution 1
from 0 to 3 m : background (solution 0).
\item the recharge composition is entered in ph.5 line at the top cells using solution 3.
\item in Rates specify the degradation rate parameters of Orgc (0,1e-11, 1e-12, 1e-12,1e-17,1e-17 respectively). 
The formula is Orgc -1.0 CH2O(NH3)0.15 1.0.
\item use two surface sites : S_s and S_w with 1.59e-5 and .0018 respective background concentrations

\begin{table}[hb]
\caption{Chemical composition for PHT3D Exercise 7}
\begin{center}
\begin{tabular}{lc} 
\toprule
Parameter  & Background & Solution 1 & Solution 2 & Solution 3\\
\midrule
C(4)      &   1e-4 &   1e-4  &   1e-4 &   1e-4\\
Cl      &   0 &   0  &  0  &   1e-4 \\
Fe(2)      &   0 &   0  &  0  &   0 \\
Fe(3)      &   0 &   0  &  0  &   0 \\
Na      &   2.3e-4 &   2.3e-4 &   2.3e-4 &   3.13e-4 \\
O(0)      &   5e-4 &   0 &   5e-4 &   5e-4\\
Orgc      &   0 &   1e-4 &   0  &  0 \\
S(6)      &   1e-4 &   1e-4  &   1e-4 &   1e-4\\
S(-2)      &   0 &   0  &  0  &   0 \\
Zn      &   0 &   8e-6  &  8e-6  &   0 \\
pH      &   6 &   6 &  6  &   5.6 \\
pe      &   14.6 &   4 &  14.6  &   14.2 \\
\bottomrule

\end{tabular}
\label{spec_table}
\end{center}
\end{table}

